\documentclass[a4paper]{article}
\usepackage{amsmath}
\usepackage{amsfonts}
\usepackage[affil-it]{authblk}
\usepackage[backend=bibtex,style=numeric]{biblatex}

\usepackage{geometry}
\geometry{margin=1.5cm, vmargin={0pt,1cm}}
\setlength{\topmargin}{-1cm}
\setlength{\paperheight}{29.7cm}
\setlength{\textheight}{25.3cm}
\setlength{\parindent}{0pt}

\begin{document}
% =================================================
\title{Numerical Analysis theoretical homework \# 2}

\author{WangHao 3220103613
  \thanks{Electronic address: \texttt{3220103613@zju.edu.cn}}}
\affil{(Information and Computing Sciences, Class 2201), Zhejiang University }


\date{Due time: \today}

\maketitle



% ============================================
\section*{I.}
For $f \in C^2[x_0,x_1]$ and $x \in (x_0,x_1)$, 
linear interpolation of $f$ at $x_0$ and $x_1$ yields

$$
f(x) - p_1(f;x) = \frac{f''(\xi(x))}{2} (x - x_0)(x-x_1).
$$

Consider the case $f(x) = \frac{1}{x}, x_0 = 1, x_1 = 2$.


\begin{itemize}
    \item Determine $\xi(x)$ explicitly.
    \item Extend the domain of $\xi$ continuously from $(x_0,x_1)$ to $[x_0,x_1].$ 
    Find $\max \xi(x)$, $\min \xi(x)$, and $\max f''(\xi(x))$.
\end{itemize}

\subsection*{Answer of I-a}

Use the Newton formula.

$$
\begin{array}{c|cc}
x_i & f[x_i] & f[x_i, x_{i+1}] \\
\hline
1 & 1 & -\frac{1}{2} \\
2 & \frac{1}{2} \\
\end{array}
$$

Then substitute them into the Newton formula, we have
\begin{align*}
  p_1(x)&=f[x_0]+f[x_0,x_1](x-x_0)\\
  &=\frac{3}{2}-\frac{1}{2}x
\end{align*}

Taking the second derivative with respect to f, we have $$f''(x)=2x^{-3}$$

Substitute them into
$$
f(x) - p_1(f;x) = \frac{f''(\xi(x))}{2} (x - x_0)(x-x_1).
$$
We have $\xi=\sqrt[3]{2x}$, where $x \in (1,2)$.

\subsection*{Answer of I-b}
As $\xi=\sqrt[3]{2x}$ is continuous, we can extend it's domain to $[1,2]$.\\
Then $\max\xi=\sqrt[3]{4}$, $\min\xi=\sqrt[3]{2}$.\\
And $f''(\xi)=\frac{1}{x}$, thus $\max f''(\xi)=1$.

% ============================================
\section*{II.}
Let $\mathbb{P}_m^+$ be the set of all polynomials of degree $\leq m$
that are non-negative on the real line,
$$
\mathbb{P}_m^+ = \{p : p \in \mathbb{P}_m, \forall x \in \mathbb{R}, p(x) \geq 0\}.
$$
Find $p \in \mathbb{P}_{2n}^+$ such that $p(x_i) = f_i$ for $i = 0,1\ldots,n$
where $f_i \geq 0$ and $x_i$ are distinct points on $\mathbb{R}$.
\subsection*{Answer of II}

The Lagrange formula:
$$p_n(x)=\sum_{i=0}^n f_i \ell_i(x)$$
Let $b_i(x)=\ell_i(x)^2$, for $i= 0,1,\ldots,n$\\
As $f_i \geq 0$, for $i= 0,1,\ldots,n$\\
We can replace $\ell_i(x)$ with $b_i(x)$, 
and the Lagrange formula still holds.

So we find
\begin{align*}
  p(x)&=\sum_{i=0}^n f_i b_i(x)\\
  &=\sum_{i=0}^n f_i [\prod_{j=0,j\neq i}^{n} \frac{x-x_j}{x_i-x_j}]^2
\end{align*}
\\
which satisfies $p(x) \in \mathbb{P}_{2n}^+$.

% ============================================
\section*{III.}
Consider $f(x) = e^x$.
\begin{itemize}
  \item Prove by induction that
  $$\forall t \in \mathbb{R},\quad f[t,t+1,\ldots,t+n]=\frac{(e-1)^n}{n!}e^t$$
  \item From Corollary 2.22 we know
  $$ \exists \xi \in (0,n) \quad s.t. \quad f[0,1,\ldots,n]=\frac{1}{n!}f^{(n)}(\xi).$$
  Determine $\xi$ from the above two equations. Is $\xi$
  located to the left or to the right of the midpoint $n/2$?
\end{itemize}

\subsection*{Answer of III-a}

Induction:

\subsubsection*{Base case}
For $n=0$,\\
$LHS=f[t]=f(t)=e^t$,\\
$RHS=e^t$.\\
The equation holds.

\subsubsection*{Inductive hypothesis}
Assume the equation holds for $ n=k $.

\subsubsection*{Inductive step}
For $n=k+1$,
\begin{align*}
  f[t,t+1,\ldots,t+k+1]
  &=\frac{f_k^{[t+1]}-f_k^{[t]}}{t+k+1-t}\\
  &=\frac{\frac{(e-1)^k}{k!}(e^{t+1}-e^t)}{k+1}\\
  &=\frac{(e-1)^{k+1}}{(k+1)!}e^t
\end{align*}

Thus, $\forall t, \forall n < \infty$, the equation holds.

\subsection*{Answer of III-b}
In $$f[t,t+1,\ldots,t+n]=\frac{(e-1)^n}{n!}e^t$$
Let $t=0,n=n$,
substitute it into $$f[0,1,\ldots,n]=\frac{1}{n!}f^{(n)}(\xi), \text{ where } \xi \in (0,n)$$
We can solve that $$\xi=n\ln(e-1)$$
As $\ln(e-1)=0.54\ldots$ , $\xi$ is located to the right of $n/2$.

% ============================================
\section*{IV.}
Consider $f(0) = 5$, $f(1) = 3$, $f(3) = 5$, $f(4) = 12$.
\begin{itemize}
  \item Use the Newton formula to obtain $p_3(f; x)$;
  \item The data suggest that $f$ has a minimum in
  $x \in (1,3)$. Find an approximate value for the location $x_{min}$ of the minimum.
\end{itemize}

\subsection*{Answer of IV-a}
Construct the difference quotient table.
$$
\begin{array}{c|cccc}
x_i & f[x_i] & f[x_i, x_{i+1}] & f[x_i, x_{i+1},x_{i+2}] & f[x_i, x_{i+1},x_{i+2},x_{i+3}]\\
\hline
0 & 5 & -2 & 1 &\frac{1}{4}\\
1 & 3 & 1 &2 \\
3 & 5 & 7\\
4 & 12 \\
\end{array}
$$
Substitute the difference quotient into Newton formula, \\
we obtain $p_3(f;x)=\frac{1}{4}x^3-\frac{9}{4}x+5$.

\subsection*{Answer of IV-b}
$p_3'(f;x)=\frac{3}{4}x^2-\frac{9}{4},\quad p_3''(f;x)=\frac{3}{2}x$\\
Solve $p_3'(f;x)=0$, we have $x=\pm \sqrt[]{3}$, as $p_3''(f;\sqrt[]{3})>0$,\\
we find the approximate $x_{min}=\sqrt[]{3}$ 

% ============================================
\section*{V.}
Consider $f(x) = x^7$.
\begin{itemize}
  \item Compute $f[0,1,1,1,2,2]$.
  \item We know that this divided difference is expressible 
  in terms of the 5th derivative of $f$ evaluated
  at some $\xi \in (0,2)$. Determine $\xi$.
\end{itemize}

\subsection*{Answer of V-a}
Construct the difference quotient table.
$$
\begin{array}{c|cccccc}
x_i & f[x_i] & f[x_i, x_{i+1}] & f[x_i, x_{i+1},x_{i+2}] & f[x_i, x_{i+1},x_{i+2},x_{i+3}] & f[x_i, x_{i+1},x_{i+2},x_{i+3},x_{i+4}] & f[x_i, x_{i+1},x_{i+2},x_{i+3},x_{i+4},x_{i+5}]\\
\hline
0 & 0 & 1 & 6 & 15 & 42 & 30 \\
1 & 1 & 7 & 21 & 99 & 102 \\
1 & 1 & 7 & 120 & 201\\
1 & 1 & 127 & 321\\
2 & 128 & 448 \\
2 & 128 \\
\end{array}
$$

Thus $f[0,1,1,1,2,2]=30$

\subsection*{Answer of V-b}
The equation is $$f[0,1,1,1,2,2]=\frac{f^{(5)}(\xi)}{5!}, \text{ where } \xi \in (0,2)$$
Substitute $f^{(5)}(x)=2520x^2$ into it, we can solve that
$$\xi=\sqrt[]{\frac{10}{7}}$$

% ============================================
\section*{VI.}
$f$ is a function on $[0,3]$ for which one knows that
$$f(0) = 1,f(1) = 2,f'(1) = -1,f(3) = f'(3) = 0.$$
\begin{itemize}
  \item Estimate $f(2)$ using Hermite interpolation.
  \item Estimate the maximum possible error of the above
  answer if one knows, in addition, that $f \in \mathcal{C}^5[0,3]$
  and $|f^{(5)}(x)| \leq M$ on $[0,3]$. Express the answer
  in terms of $M$.
\end{itemize}

\subsection*{Answer of VI-a}
Construct the difference quotient table.
$$
\begin{array}{c|ccccc}
x_i & f[x_i] & f[x_i, x_{i+1}] & f[x_i, x_{i+1},x_{i+2}] & f[x_i, x_{i+1},x_{i+2},x_{i+3}] & f[x_i, x_{i+1},x_{i+2},x_{i+3},x_{i+4}] \\
\hline
0 & 1 & 1 & -2 & \frac{2}{3} & -\frac{5}{36} \\
1 & 2 & -1 & 0 & \frac{1}{4} \\
1 & 2 & -1 & \frac{1}{2} \\
3 & 0 & 0 \\
3 & 0 \\
\end{array}
$$

So we have $$p_4(x)=1+x-2x(x-1)+\frac{2}{3}x(x-1)^{2}-\frac{5}{36}x(x-1)^{2}(x-3)$$ 
and $$f(2)\approx p_4(2)=\frac{11}{18}$$

\subsection*{Answer of VI-b}
$$R(x)=\frac{f^{(5)}(\xi)}{5!}\prod_{i=0}^{4}(x-x_i), \text{ where } x=2$$

Thus $\max R(2)=\frac{M}{60} $


% ============================================
\section*{VII.}
Define forward difference by
$$\Delta f (x) = f (x + h) - f (x),$$
$$\Delta^{k+1} f (x) = \Delta \Delta^k f (x) = \Delta^k f (x + h) - \Delta^k f (x)$$
and backward difference by
$$\nabla f (x) = f (x) - f (x - h),$$
$$\nabla^{k+1} f (x) = \nabla \nabla^k f (x) = \nabla^kf (x) - \nabla^k f (x - h).$$
For $x_j = x + jh$, prove
$$\Delta^k f (x) = k! h^k f[x_0, x_1,\ldots , x_k],$$
$$\nabla^k f (x) = k! h^kf [x_0, x_{-1},\ldots, x_{-k}].$$

\subsection*{Answer of VII-a}

\subsubsection*{Base case}
For $k=1$,\\
$LHS=\Delta f(x_0)=f(x_1)-f(x_0)$,\\
$RHS=hf[x_0,x_1]=h\frac{f[x_1]-f[x_0]}{x_1-x_0}=f(x_1)-f(x_0)$.\\
The equation holds.

\subsubsection*{Inductive hypothesis}
Assume the equation holds for $ k=n $.

\subsubsection*{Inductive step}
For $k=n+1$,
\begin{align*}
  \Delta^{n+1}f(x)&=\Delta^nf(x_1)-\Delta^nf(x_0)\\
  &=n!h^n(f[x_1,x_2,\ldots,x_{n+1}]-f[x_0,x_1,\ldots,x_n])\\
  &=(n+1)!h^{n+1}f[x_0,x_1,\ldots,x_{n+1}]
\end{align*}

Thus, the equation holds.

For backward difference, the induction is similar.


% ============================================
\section*{VIII.}
Assume $f$ is differentiable at $x_0$. Prove
$$\frac{\partial}{\partial x_0}f[x_0,x_1,\ldots,x_n]=f[x_0,x_0,x_1,\ldots,x_n] $$
What about the partial derivative with respect to one
of the other variables?

\subsection*{Answer of VIII}
\begin{align*}
  LHS&=\lim_{\Delta x \to 0} \frac{f[x_0+\Delta x,x_1,\ldots,x_n]-f[x_0,x_1,\ldots,x_n]}{x_0+\Delta x-x_0}\\
  &=\lim_{\Delta x \to 0}f[x_0,x_1,\ldots,x_n,x_0+\Delta x]
\end{align*}

\begin{align*}
  Rhs&=\frac{f[x_0,x_1,\ldots,x_n]-f[x_0,x_0]}{x_n-x_0}\\
  &=\frac{f[x_0,x_1,\ldots,x_n]-f'(x_0)}{x_n-x_0}
\end{align*}
Thus $RHS$ is continuous with respect to $x_0$ when $x_0 \neq x_i, i=1,2,\ldots,n $

So
\begin{align*}
  Rhs&=f[x_0,x_1,\ldots,x_n,x_0]\\
  &=f[x_0,x_1,\ldots,x_n,\lim_{\Delta x \to 0}(x_0+\Delta x)]\\
  \text{(continuous)}&=\lim_{\Delta x \to 0}f[x_0,x_1,\ldots,x_n,x_0+\Delta x]\\
  &=LHS
\end{align*}

% ============================================
\section*{IX.}
min-max problem. For $n \in \mathbb{N}^+$ and a fixed $a_0 \neq 0$,
determine $\min \max_{x\in [a,b]} |a_0x^n + a_1x^{n-1} + \ldots + a_n|$
where the minimum is taken over all $a_i \in \mathbb{R}, i =1, 2,\ldots , n$.

\subsection*{Answer of IX}

$\min_{a_i} \max_{x\in [a,b]} |a_0x^n + a_1x^{n-1} + \ldots + a_n|
=min_{a'_i} \max_{x\in [a,b]} |a'_0(\frac{2}{b-a})^n(x-\frac{a+b}{2})^n + a'_1(\frac{2}{b-a})^{n-1}(x-\frac{a+b}{2})^{n-1} + \ldots + a'_n|$
(Because of the randomicity of $a_i$, we can solve $a'_i$ with the order of $a'_0$ to $a'_n$, and get the original formula.)\\
And $a'_0 = (\frac{2}{b-a})^n a_0$\\
Then make transform operation $$x=\frac{a+b}{2}+\frac{b-a}{2}t, \text{ where } t \in [-1,1]$$
\begin{align*}
  \text{the original formula}&=\min_{a'_i} \max_{t \in [-1,1]} |a'_0t^n + a'_1t^{n-1} + \ldots + a'_n|\\
  &=a'_0\min_{a'_i} \max_{t \in [-1,1]} |t^n + a'_1t^{n-1} + \ldots + a'_n|\\
  &=a'_0\min_{p_n(t)} \max_{t \in [-1,1]} |p_n(t)| \text{, in which } p_n(t) \in \tilde{P_n}
\end{align*}

By Theorem 2.47,
\begin{align*}
  \text{the original formula}&=a'_0 \frac{1}{2^{n-1}}\\
  &=(\frac{2}{b-a})^n a_0 \frac{1}{2^{n-1}}\\
  &= 2(\frac{b-a}{4})^n a_0
\end{align*}
When $p_n(t)=\frac{T_n(t)}{2^{n-1}}$, equality holds.

% ============================================
\section*{X.}
Imitate the proof of Chebyshev Theorem.\\
Express the Chebyshev polynomial of degree $n \in \mathbb{N}$ as
a polynomial $T_n$ and change its domain from $[-1, 1]$ to
$\mathbb{R}$. For a fixed $a > 1$, define $\mathbb{P}_n^a
 := \{p \in \mathbb{P}_n : p(a) = 1\}$
and a polynomial $\hat{p}_n(x) \in \mathbb{P}_n^a$,
$$\hat{p}_n(x) := \frac{T_n(x)}{T_n(a)}.$$
For $||f||_\infty = \max_{x\in [-1,1]} |f (x)|$, prove
$$\forall p \in \mathbb{P}_n^a, \quad ||\hat{p}_n||_\infty \leq ||p||_\infty.$$

\subsection*{Answer of X}
Assume $\exists p $ s.t. 
$$||p||_\infty < ||\hat{p}_n||_\infty$$
i.e.
$$ \max_{x\in [-1,1]} |p (x)| < \frac{1}{T_n(a)}$$
Consider $Q(x)=\frac{T_n(x)}{T_n(a)}-p(x) \quad \in \mathbb{P}_n $\\
$Q(x'_k)=\frac{(-1)^k}{T_n(a)}-p(x)$, the sign varies from $n+1$ points.\\
So it have $n$ roots between $[-1,1]$.\\
And we can notice $Q(a)=0$\\
So Q have $n+1$ roots, we obtain a contradiction, thus the proof is complete.

% ============================================
\section*{XI.}
Give a detailed proof of Lemma 2.53.

\subsection*{Answer of XI}
\begin{align*}
  LHS&=\binom{n-1}{k}t^k(1-t)^{n-1-k}\\
  &=\frac{k+1}{n}\binom{n}{k+1}t^k(1-t)^{n-1-k}\\
  &=\frac{k+1}{n}\binom{n}{k+1}t^k(1-t)^{n-1-k}[(1-t)+t]\\
  &=\frac{k+1}{n}\frac{n-k}{k+1}\binom{n}{k}t^k(1-t)^{n-k}+\frac{k+1}{n}\binom{n}{k+1}t^{k+1}(1-t)^{n-1-k}\\
  &=RHS
\end{align*}

% ============================================
\section*{XII.}
Give a detailed proof of Lemma 2.55.

\subsection*{Answer of XII}
$\forall k \geq 1$, we have
\begin{align*}
  LHS&=\int_{0}^{1}\binom{n}{k}t^{k}(1-t)^{n-k}dt\\
  &=\binom{n}{k}(\left[ -t^k \cdot \frac{(1-t)^{n-k+1}}{n-k+1} \right]_0^1-\int_0^1 k t^{k-1} \cdot \frac{(1-t)^{n-k+1}}{n-k+1} \,dt)\\
  &=\frac{k}{n-k+1}\int_0^1 \binom{n}{k}t^{k-1}(1-t)^{n-k+1}dt\\
\end{align*}
Using this equation repeatedly, we have
\begin{align*}
LHS&=\int_0^1 t^0 (1-t)^n \, dt \\
&= \int_0^1 (1-t)^n \, dt = \frac{1}{n+1}
\end{align*}

\end{document}